\documentclass{article}
\usepackage[T1]{fontenc}
\usepackage{lmodern}
\usepackage{graphicx}
\usepackage{amsmath,amssymb}
\usepackage[a4paper,margin=2cm]{geometry}
\usepackage[absolute,overlay]{textpos}
\setlength{\TPHorizModule}{1mm}
\setlength{\TPVertModule}{1mm}
\usepackage[indonesian]{babel}
\usepackage{hyperref}
\setlength{\emergencystretch}{2em}
% unit: Halaman 33

\begin{document}

% === Halaman 33 (python/native) ===

• Contoh: Misalkan P(2, 4) dan Q(5, 9) adalah dua buah titik pada kurva eliptik 

y2  x3 + x + 6  (mod 11). Tentukan P  + Q dan 2P.


 Jawab: 

    (a) P + Q = R

  m   (9 – 4)/(5 – 2) (mod 11) =  5/3  (mod 11) = 5  3–1 (mod 11)



                      = 5  4 (mod 11)  9 (mod 11) 

 P + Q = R, koordinat Titik R: 

     xr  m2 – xp – xq (mod 11)  81 – 2 – 5 mod 11  8 (mod 11)

     yr  m(xp – xr) – yp (mod 11)  9(2 – 8) – 4 mod 11  -58 (mod 11) 







                 8 (mod 11)

 Jadi,  R(8, 8) 

33

\end{document}
